\section{Confluence of \texorpdfstring{$\lbox$}{lambda-box}}
\label{sec:confluence}

We prove \cref{thm:confluence} by the Tait--Martin-L{\"o}f method in
Takahashi's formulation \cite{Takahashi1995}, extended to the $\iota$,
$\mu$ and $\delta$ rules.  The proof is that of the companion paper,
reproduced for self-containedness; it depends on the constructor
alphabet only through arities, so the indexed extension changes
nothing.  Recall that the $\delta$-rule unfolds a constant to the
erasure of its body, $c \lred \erase{t_c}$, and that bodies mention only
earlier constants; no property of this ordering is needed beyond the
fact that each constant has a \emph{fixed} unfolding.

\begin{definition}[Parallel reduction]
\label{def:parred}
The relation $\Rightarrow$ on $\lbox$ terms is defined inductively:
\begin{enumerate}[leftmargin=2em,itemsep=0.2ex]
\item $x \Rightarrow x$, \quad $\bx \Rightarrow \bx$, \quad $c
  \Rightarrow c$;
\item $c \Rightarrow \erase{t_c}$ for a defined constant $c$;
\item $\lambda x.e \Rightarrow \lambda x.e'$ if $e \Rightarrow e'$;
\item $e_1\,e_2 \Rightarrow e_1'\,e_2'$ if $e_i \Rightarrow e_i'$;
\item $\wcns{c}(\tup{e}) \Rightarrow \wcns{c}(\tup{e}')$ if $e_i
  \Rightarrow e_i'$ componentwise;
\item $\mathsf{case}(e; \tup{br}) \Rightarrow \mathsf{case}(e';
  \tup{br}')$ if $e \Rightarrow e'$ and branch bodies reduce
  componentwise;
\item $\mathsf{fix}\,f.e \Rightarrow \mathsf{fix}\,f.e'$ if $e
  \Rightarrow e'$;
\item $(\lambda x.e)\,a \Rightarrow \subst{e'}{a'}{x}$ if $e \Rightarrow
  e'$ and $a \Rightarrow a'$;
\item $\mathsf{case}(\wcns{c}_i(\tup{e}); \{\wcns{c}_j(\tup{x}_j)
  \Rightarrow b_j\}_j) \Rightarrow \subst{b_i'}{\tup{e}'}{\tup{x}_i}$ if
  $e_l \Rightarrow e_l'$ and $b_i \Rightarrow b_i'$;
\item $\mathsf{fix}\,f.e \Rightarrow
  \subst{e'}{\mathsf{fix}\,f.e'}{f}$ if $e \Rightarrow e'$.
\end{enumerate}
\end{definition}

\begin{lemma}
\label{lem:parred-basic}
(i) $e \Rightarrow e$ for all $e$.\quad
(ii) $\lred\; \subseteq\; \Rightarrow\; \subseteq\; \lreds$.\quad
(iii) If $e \Rightarrow e'$ and $u \Rightarrow u'$ then
$\subst{e}{u}{x} \Rightarrow \subst{e'}{u'}{x}$.
\end{lemma}

\begin{proof}
(i) Induction on $e$ using clauses (1) and (3)--(7).  (ii) For the first
inclusion, a one-step reduction is a root rule applied inside a context:
the corresponding root clause (2), (8), (9) or (10) with reflexive
subderivations, propagated by the congruence clauses.  For the second,
induction on $e \Rightarrow e'$: each clause is realised by finitely
many $\lred$-steps.  (iii) Induction on the derivation of $e
\Rightarrow e'$.  The variable clause $x \Rightarrow x$ becomes $u
\Rightarrow u'$; clause (2) is unaffected because $\erase{t_c}$ is
closed; the congruence clauses follow from the induction hypothesis;
the redex clauses use the standard substitution identities, e.g.\
$\subst{(\subst{e'}{a'}{y})}{u'}{x} =
\subst{(\subst{e'}{u'}{x})}{\subst{a'}{u'}{x}}{y}$ for $y \notin
\fv{u'} \cup \{x\}$.
\end{proof}

\begin{definition}[Complete development]
\label{def:development}
The complete development $e^{\ast}$ is defined by recursion on $e$:
\begin{align*}
  x^{\ast} &= x, \qquad \bx^{\ast} = \bx, \qquad
  c^{\ast} = \erase{t_c},\\
  (\lambda x.e)^{\ast} &= \lambda x.e^{\ast}, \qquad
  \wcns{c}(\tup{e})^{\ast} = \wcns{c}(\tup{e}^{\ast}), \qquad
  (\mathsf{fix}\,f.e)^{\ast} =
    \subst{e^{\ast}}{\mathsf{fix}\,f.e^{\ast}}{f},\\
  ((\lambda x.e)\,a)^{\ast} &= \subst{e^{\ast}}{a^{\ast}}{x}, \qquad
  (e_1\,e_2)^{\ast} = e_1^{\ast}\,e_2^{\ast} \quad (e_1 \text{ not a }
    \lambda),\\
  \mathsf{case}(\wcns{c}_i(\tup{e}); \tup{br})^{\ast} &=
    \subst{b_i^{\ast}}{\tup{e}^{\ast}}{\tup{x}_i},\\
  \mathsf{case}(e; \tup{br})^{\ast} &=
    \mathsf{case}(e^{\ast}; \tup{br}^{\ast}) \quad (e \text{ not a
    constructor form}).
\end{align*}
(For a constant, the development unfolds once and does not recurse into
the unfolding; for a fix, the root $\mu$-redex is contracted once, its
residuals in the contractum being new redexes that are not developed.)
\end{definition}

\begin{lemma}[Triangle property]
\label{lem:triangle}
If $e \Rightarrow e'$ then $e' \Rightarrow e^{\ast}$.
\end{lemma}

\begin{proof}
Induction on the derivation of $e \Rightarrow e'$, with a case analysis
on the last clause used.
\begin{itemize}[leftmargin=2em]
\item Clause (1): for variables and $\bx$ immediate; for $c \Rightarrow
  c$ we need $c \Rightarrow c^{\ast} = \erase{t_c}$, which is clause
  (2).
\item Clause (2): $e' = \erase{t_c} = c^{\ast}$, and $e' \Rightarrow e'$
  by \cref{lem:parred-basic}(i).
\item Clauses (3), (5), (7), and clause (6) with non-constructor
  scrutinee: by the induction hypothesis componentwise.  Two
  subtleties: for $\mathsf{fix}\,f.e \Rightarrow \mathsf{fix}\,f.e'$
  (clause (7)) the target is $(\mathsf{fix}\,f.e)^{\ast} =
  \subst{e^{\ast}}{\mathsf{fix}\,f.e^{\ast}}{f}$, reached from
  $\mathsf{fix}\,f.e'$ by clause (10) with $e' \Rightarrow e^{\ast}$
  (induction hypothesis).  For $\mathsf{case}(e; \tup{br})$ with $e$
  not a constructor form, the scrutinee's reduct $e''$ may or may not
  be a constructor form, but the development
  $\mathsf{case}(e^{\ast}; \tup{br}^{\ast})$ is reached by clause (6)
  either way.
\item Clause (4) with $e_1 = \lambda x.b$: then $e_1 \Rightarrow e_1'$
  is by clause (3), $e_1' = \lambda x.b'$, and $e^{\ast} =
  \subst{b^{\ast}}{a^{\ast}}{x}$; from $e' = (\lambda x.b')\,a'$ apply
  clause (8) with $b' \Rightarrow b^{\ast}$, $a' \Rightarrow a^{\ast}$
  (induction hypotheses).
\item Clause (4) with $e_1$ not a $\lambda$: $e^{\ast} =
  e_1^{\ast}\,e_2^{\ast}$; apply clause (4) with the induction
  hypotheses.
\item Clause (8): $e = (\lambda x.b)\,a$, $e' = \subst{b'}{a'}{x}$; by
  induction hypotheses $b' \Rightarrow b^{\ast}$ and $a' \Rightarrow
  a^{\ast}$, so \cref{lem:parred-basic}(iii) gives $e' \Rightarrow
  \subst{b^{\ast}}{a^{\ast}}{x} = e^{\ast}$.
\item Clause (9): as clause (8), with the componentwise induction
  hypotheses for the constructor arguments and the branch, and
  \cref{lem:parred-basic}(iii) extended to simultaneous substitutions.
  The scrutinee $\wcns{c}_i(\tup{e})$ can only have been reduced
  componentwise (clause (5)), so the case analysis in
  \cref{def:development} is stable.
\item Clause (10): $e = \mathsf{fix}\,f.b$, $e' =
  \subst{b'}{\mathsf{fix}\,f.b'}{f}$; by the induction hypothesis $b'
  \Rightarrow b^{\ast}$, whence $\mathsf{fix}\,f.b' \Rightarrow
  \mathsf{fix}\,f.b^{\ast}$ by clause (7), and
  \cref{lem:parred-basic}(iii) gives $e' \Rightarrow
  \subst{b^{\ast}}{\mathsf{fix}\,f.b^{\ast}}{f} = e^{\ast}$.
\end{itemize}
\end{proof}

\begin{proof}[Proof of \cref{thm:confluence}]
By \cref{lem:triangle}, $\Rightarrow$ has the diamond property: if $e
\Rightarrow e_1$ and $e \Rightarrow e_2$ then $e_1 \Rightarrow e^{\ast}$
and $e_2 \Rightarrow e^{\ast}$.  A relation with the diamond property
has confluent reflexive--transitive closure (tiling), and by
\cref{lem:parred-basic}(ii) the closure of $\Rightarrow$ is $\lreds$.
The characterisation of $\lconv$ by common reducts follows by induction
on the conversion zig-zag (Church--Rosser).
\end{proof}
